\chapter[Complexity Without Scriptures]{Does a Tradition Without Scriptures Lack Complexity?}
\section{A Drum in a Display Case}
First, pick up an object. This time, however, it is not in your hands but behind glass.

Imagine entering the basement of a large museum. The lights are dim, and a wooden mask lies in a display case. Slightly wider than your face, it is carved from driftwood and painted in faded red and white. Two or three wooden rings encircle the face. Feathers once stood in them; now only a few broken shafts remain. The label says, roughly: mask, southwestern Alaska, Yup'ik people, collected in the nineteenth century, used in shamanic ceremonies.

You look at it through the glass; it looks back at you.

Open a typical introduction to world religions and you probably will not find the Yup'ik. There will be whole chapters on Buddhism, Christianity, and Islam, perhaps a page or two on "other traditions." Alaska's Indigenous peoples, Aboriginal Australians, and peoples of the Amazon rainforest are usually compressed into one term: "primitive religion," or slightly more politely, "tribal beliefs" or "Indigenous religions." The implication is clear: these are religion's childhood, simple and crude, not yet grown into a mature form with scriptures, theology, and churches.

That implication is precisely what this chapter will examine.

The Yup'ik have no religion written down in a sacred book. They have no Bible or Quran, no seminary, no pope. But does a tradition without scripture therefore amount to a simple tradition? Is the world behind this mask really shallower than the world of a three-thousand-year-old sacred book?

Before answering, let us leave the gallery for the place where the mask truly lived.

\section{A Gathering House on a Long Winter Night}
Time: a winter night more than a century ago. Place: tundra between the Yukon and Kuskokwim rivers in southwestern Alaska, in a gathering house partly sunk into the ground, called a qasgiq in Yup'ik. Outside, the dark air is thirty or forty degrees Celsius below zero, and wind drives snow sideways. Inside, a few seal-oil lamps cast flickering light on wooden walls.

The house is crowded. Men have bare upper bodies because fire and body heat make the air hot; women and children gather toward the back. The warmth carries the smells of seal oil, damp wood, and dozens of sweating bodies. Such a gathering house is the village's heart: men work, eat, and discuss hunting here. Tonight the whole village has assembled for a winter festival.

Drums begin, not one but many. The drummers sing about a whale that long ago agreed to a hunter's request. Two masked dancers step from darkness into the light. Wooden rings and feathers turn in the firelight; the masks' human faces have half-closed eyes and ambiguous smiles. As the rhythm changes, so do the gestures: a turn of the wrist imitates a loon's wings; a sinking body depicts an animal struck by a harpoon.

The elders seated in a corner can name every gesture, identify the family to which each melody belongs, and explain which encounter a mask recounts: between a person and a seal, or a person and the spirit of the wind. They know who sang the song last year, who should inherit it next year, which passages women may learn, and which may be voiced only on the eve of the ceremony. After the festival, some masks will remain where they are until they decay. A mask is not an artwork but a temporary dwelling for a particular encounter; once the business is finished, the dwelling should depart.

This is no "entertainment evening for primitive people." It is the public performance of a precisely functioning knowledge system: how animals make agreements with humans, where people go after death, the particular temperaments of wind, rivers, and sea, and the kindnesses a family has owed and repaid over centuries. Not a single line is written on paper. Every "text" lives in hundreds of bodies, voices, gestures, and drumbeats.

Remember this hot, crowded room. The rest of the chapter answers the question it poses.

\section[No Sacred Books, No Thought?]{Does Having No Sacred Books Mean Having No Thought?}
Let us now set out the tension, without rushing to an answer.

On one side stands a ranking we have absorbed almost since childhood. Open a chart of world religions and you see major trunks: Judaism, Christianity, Islam, Hinduism, Buddhism. Each corresponds to scriptures, founders, theology, and organization. In the margins cluster a few gray entries: animism, shamanism, ancestor worship. The ranking seems self-evident: literate religions are advanced; oral, local traditions without written scripture are lower forms, as tadpoles are to frogs or drafts to finished texts.

This ranking has practical conveniences, whose case we will later state fully. On the other side, however, stand scenes like the Yup'ik gathering house. There are no tadpoles here. Safeguards for transmitting songs are stricter than many printers' proofreading procedures; elders immediately hear a mistaken word. Thought about relations between humans and animals extends to the precise ways a hunted animal's soul might return. A festival calendar fits closely with stars, fish runs, and bird migrations. How is this simple?

The real question emerges:

\textbf{Might the ruler we use to measure complexity itself be crooked?}

If its markings ask whether there are written scriptures, full-time clergy, and systematic theology, traditions without writing automatically score zero. Yet perhaps this ruler measures not complexity but resemblance to religions familiar to us. A culture that does not preserve knowledge in writing does not therefore possess less knowledge, any more than someone who remembers everything without a notebook has no memory. The problem is visibility: knowledge on paper can be seen; knowledge in bodies and voices cannot be seen through museum glass.

Dig deeper and a more uncomfortable issue appears. These traditions left no official documents of their own. Almost everything we now know about them comes through outsiders' pens: missionaries, merchants, anthropologists, and colonial officials. We therefore see not the gathering house itself but someone else's photograph taken through glass. The photographer's position, lens, and choice of moment all shape what we see. Between an elder's entire known cosmos and three pages written by a missionary who visited for three summers lies the distance this chapter repeatedly measures.

Before continuing, answer for yourself: which do you trust more, the gray lines labeled "primitive religion" in the textbook chart or the steaming gathering house? Where does your first reaction come from?

\section{Who Records, Who Is Recorded?}
At least two positions clash over these traditions, on a profoundly unequal field: only one side has a printing press.

\textbf{Voice one: the outside recorder.} Nineteenth-century missionaries, colonial officials, and early anthropologists wrote most of the accounts we can read today. In their descriptions, these traditions often did not even deserve the name religion. Many missionary reports said, in effect, "These people have no religion, only superstition and tales of ghosts." This reveals less about local life than about the recorder's mold. Their definition of religion was translated from Christianity: one God, scripture, a church, a catechism. Measure Yup'ik gathering houses, Australian songlines, or Yoruba divination poetry with that mold and nothing fits, producing the conclusion that religion is absent. It is like using a butterfly net to declare the sea lifeless.

\textbf{Voice two: Indigenous voices.} Indigenous voices did survive, but by more winding paths. In the 1920s, the Danish explorer Knud Rasmussen crossed the Arctic, recording many Inuit songs and oral accounts. Having grown up in Greenland and speaking its language, he was unusual among recorders. In the same period, anthropologist Franz Boas and his Kwakwaka'wakw collaborator George Hunt worked together over decades, leaving an immense body of primary texts from the Indigenous peoples of North America's Northwest Coast, many written by Hunt himself in his community's language. Here Indigenous narrators finally approached speaking for themselves. But note both qualifications: finally and approached. Even the best records were selected: what to record, what to omit, and what narrators themselves judged outsiders should not hear. Much sacred knowledge was precisely the material reserved for authorized people within the community, so narrators had reason to withhold it. What remains on paper is never the entire tradition, only the intersection of what it was willing to show and what recorders were willing to copy.

Now practice stating the other side's case fully. Our subjects are the nineteenth-century missionary recorders, whom it is easy to dismiss wholesale.

Present their circumstances as fairly as possible. First, many genuinely spent years learning local languages and produced the earliest dictionaries and grammars. Without them, many languages' later prospects for preservation would have been worse. Second, they encountered something for which their education had not prepared them: a worldview weaving people, animals, winds, and ancestors into one web. Within their only available vocabulary, "superstition" was indeed the nearest box. Third, a record is better than none. When Indigenous communities rebuild traditions today, the available materials are often those same prejudiced notes, because so many elders died amid disease and forced assimilation.

All three points are true. Yet they do not repair the fundamental flaw: someone holding both the power to explain you and the power to transform you cannot produce a neutral record. The period when missionaries wrote "there is no religion here" was also the period when churches and governments sent Indigenous children to boarding schools, banned their languages, and burned ceremonial objects. Every "superstition" in a report could license the next prohibition. Understanding recorders does not remove the need to question their records; both tasks must proceed together.

\section[Thinking Bridge: Source Checkpoints]{Thinking Bridge: How Many Checkpoints Does Information Pass?}
When we have only photographs and cannot visit the scene, we need a portable tool. Historians call it source criticism. Think of parcel tracking: information passes through several checkpoints on its way to you, and at each it may be opened, repackaged, or damaged. When reading an account of someone else's tradition, ask about each checkpoint.

\textbf{Checkpoint one: who is speaking?} Is the narrator a community elder or a passing merchant? Are they entitled within their community to tell this material? Strict rules often determine who may tell which stories. Outsiders may therefore hear a version intended for public release.

\textbf{Checkpoint two: who is recording?} Does the recorder understand the language, and why have they come? A missionary report, a colonial archive, an anthropologist's field notes, and a travel program produce four different accounts of the same ceremony.

\textbf{Checkpoint three: who is translating?} Crossing languages inevitably loses something. "Dreamtime," used for Aboriginal Australian traditions, is a famous translation mishap. Late-nineteenth-century anthropologists used it to render an Arrernte term whose reference was not dreaming but a whole interweaving of creative order, laws of the land, and present reality. Yet "Dreamtime" made a serious philosophy of land sound like a bedtime story. This mistranslation has circulated for over a century, continuing to diminish how intelligent these traditions appear.

\textbf{Checkpoint four: who decides what is published?} Recording does not guarantee printing. The texts selected, introductions attached, and classification labels applied shape every later reader's first impression. Readers approach the same unchanged words differently in a series called "Myths and Legends" than in one called "Philosophical Classics."

\textbf{Checkpoint five: who has been cut out?} Groups without surviving records are silent in the archive. Women, children, and lower-status families often never had the chance to get their versions through even the first checkpoint.

Next time you encounter an assertion about any community's traditions---what a people believes or where a religion originated---neither accept nor reject it immediately. Send it through these five checkpoints. The purpose is not to throw every old record away but to establish what a source can and cannot prove. A missionary report might establish that a particular ceremony took place somewhere in a particular month of 1903. It can hardly establish that the people involved lacked complex thought. Distinguishing these is the tool's entire point.

\section{The Authenticity Puzzle of a Famous Speech}
Now use your new tool on a short "primary source" you may have encountered.

Perhaps the world's most widely circulated "Indigenous voice" is the speech called Chief Seattle's Declaration. Seattle, or Si'ahl, was a nineteenth-century leader of the Duwamish and Suquamish peoples on America's Northwest Coast. In the circulating version, he spoke in 1854 to white officials seeking to purchase land. A familiar Chinese rendering includes this passage:

\begin{quote}
"The land does not belong to humanity; humanity belongs to the land. All things are connected, as blood ties join a family. Humanity did not weave life's web; it is only one thread within it. Whatever it does to that web, it does to itself."
\end{quote}
The passage is strikingly beautiful. Printed in environmental textbooks, hung in offices, and turned into posters, it has almost become the global spokesperson for "Indigenous wisdom."

Now send it through the five checkpoints. The result is chilling.

There was indeed a meeting about land in 1854, and Chief Seattle spoke. But no shorthand record was made. The sole white recorder, Henry Smith, claimed to reconstruct it from notes, yet published his version only in the 1880s, more than thirty years later, acknowledging his own literary embellishment. How many checkpoints have already failed? There is more: most famous lines quoted today do not appear even in Smith's questionable version. They come from a script written in 1972 by a playwright named Ted Perry for an environmental film. Perry later repeatedly acknowledged publicly that these were his words, not Seattle's. The globally quoted "Indigenous maxim" about humanity not weaving life's web was created by a twentieth-century screenwriter.

The "Indigenous voice" celebrated worldwide for half a century is thus a parcel with scarcely any original contents left after five checkpoints: Indigenous speech $\rightarrow$ questionable recollection $\rightarrow$ delayed publication $\rightarrow$ fictional screenplay $\rightarrow$ global republication. This matters not because it proves all records false but because it precisely demonstrates each checkpoint's operation. Notice an especially sharp detail: the fictional version conquered the world precisely because it perfectly matched outsiders' image of the noble Indigenous person. People often circulate not evidence but echoes of their own wishes.

This does not make the idea that humans and nature are connected false. On the contrary, such relational worldviews genuinely exist in countless Indigenous traditions, the subject of this chapter throughout. What is false is the attribution, not necessarily the idea. Separating these already makes you more discerning than most people who shared that poster.

\section{One Blank Page, Three Explanations}
First, fill the supposedly blank page with concrete examples. Claims of simple traditions become difficult to sustain once the materials are visible.

Look at West Africa. Yoruba people around Nigeria and Togo maintain the Ifa divination tradition, whose core is an enormous oral poetry archive organized into 256 units. Each contains dozens or hundreds of verses, tens of thousands altogether. Its specialist priests begin in youth and spend more than a decade memorizing them, including their myths, remedies, genealogies, and precedents. This oral library has been inscribed on UNESCO's intangible cultural heritage list. Imagine the work: memorizing an encyclopedia word for word, then retrieving passages on demand during ceremonies. No writing? Yes. No complexity? That conclusion requires refusing to count.

Look at Oceania. Maori whakapapa, genealogy, is no hobbyist's family-history booklet but an oral structure stitching people to mountains, rivers, ancestors, and land. Reciting one's genealogy means proceeding from creation to one's mother, declaring along the way which mountain, river, and ancestral canoe one belongs to. For Maori, this is not metaphor but an identity document with legal weight. In nineteenth-century New Zealand land courts, genealogical recitation served as evidence of land rights.

Look at the Andes. Quechua-speaking communities organize their world around the ayllu, a word only approximately translated as community. It includes not just living people but ancestors, land, mountain beings called apu, and Pachamama, Mother Earth. A field's poor harvest is therefore not merely a weather event but a problem in some relationship within this kinship web. Offering drink to the earth or sacrifice to a mountain fulfills obligations to kin, rather than simply expressing superstition. You would not ask someone visiting a grandmother's grave whether spirit money really transfers into an underworld account, because the point was never merely the transfer.

Look at East Asia. The Ainu of northern Japan regard bears as kamuy, divine beings visiting humans in animal form. The solemn iomante, or bear-sending ceremony, respectfully returns the spirit of a bear raised by people to the divine world, through hundreds of ritual details. Our own traditions also contain deep structures not dependent on sacred books. The Analects, "Eight Rows of Dancers," records Confucius sacrificing as if ancestors were present and honoring spirits as if spirits were present. Notice the subtlety: he does not argue whether ancestors truly attend but emphasizes "as if." Ritual meaning does not depend entirely on whether someone receives what is offered. Structurally, this closely resembles Andean offerings of drink to the earth.

The examples are before us. Now answer the central question: why are these traditions absent from world-religion charts or compressed into small gray lettering? The same blank page admits at least three nonequivalent explanations.

\textbf{Explanation one: bias in the medium, a technological explanation.} Written materials are durable, reproducible, and transportable, naturally occupying libraries and textbooks. Oral knowledge lives in bodies; when a person dies, a library burns. The chart's blank page may merely reflect an empty archival shelf, measuring the medium rather than thought. This is clear and powerful, but too gentle: it suggests a technical accident that better recording could solve. Why, then, did nobody fill the gap for centuries?

\textbf{Explanation two: someone else made the definition, a conceptual explanation.} The word religion was tailored to certain traditions: a founder, scripture, doctrine, organization, each checked off. Traditions without those components did not fail the examination; they were denied entry. This fills the first account's gap, explaining why the blanks are so systematic. But it makes the problem resemble a vocabulary misunderstanding, as if a broader definition would solve everything. The institutions defining religion were also burning masks and banning languages. This was not merely poor wording but wrongdoing through words.

\textbf{Explanation three: the power to record, a political explanation.} Who may define, record, and print records as textbooks has always been a question of power. Colonial rule needed to portray its subjects as people without real religion, because that diagnosis licensed treatment: schools, churches, and land legislation all used it as supporting documentation. The blank was engineered, not overlooked. This is the sharpest explanation, accounting for systematic patterns the others cannot. Its limitation is that treating everything as a power project obscures recorders' genuine honesty and curiosity---people like Rasmussen did exist---and turns Indigenous people into pure victims. It erases their active preservation, concealment, negotiation, and invention of new ceremonies. They were never merely recorded objects; they also decided which face to show outsiders.

Each explanation holds part of the truth. A responsible reader layers them: the medium explains why gaps are hard to fill, definitions why the gaps are uniform, and power why those gaps happen to serve conquest.

\section{Further Challenge: Two Overused Words}
Now bring out the chapter's theoretical equipment. We will dismantle two familiar words: animism and shaman. Their histories form a miniature history of colonial knowledge.

First, animism. In 1871, the British anthropologist Edward Tylor published Primitive Culture, offering a famous minimum definition of religion: broadly, belief in spiritual beings. This sounds inclusive, lowering the threshold so everyone has religion. But Tylor's tolerance placed everyone in a queue. He arranged human belief in an ascending sequence: animism, belief in spirits throughout the world, at the bottom; polytheism above it; then monotheism; and science at the summit. In this picture, the Yup'ik hunter who believes a mountain has a spirit stands at the foot of the ladder, a fossil from human thought's childhood.

This ladder dominated scholarship for decades before being dismantled step by step in the twentieth century. Critics argued that Tylor treated the European soul as a standard component and searched the world for matching parts. More fundamentally, each rung corresponded exactly to contemporary European rankings of civilization. Anthropology's ladder was a side view of the colonial map.

A major theoretical response came late in the twentieth century. After studying a foraging community in southern India, anthropologist Nurit Bird-David argued that understanding animism as a mistaken belief---wrongly supposing mountains have spirits---was itself mistaken. For these communities it is less a theory about the world than a way of relating to it. The central issue is not whether the proposition "the mountain has a spirit" is true but whether the mountain is a "someone" with whom one interacts and to whom one owes reciprocity. This is an ontology that gives priority to relationships: who comes before what. French anthropologist Philippe Descola went further, organizing ways of understanding the world into four basic types: animism, with shared interiority but differing physicality; totemism, with both shared; analogism, with correspondences between different beings; and naturalism, with shared physical nature but differing interiors. Notice the fourth: naturalism is our default scientific worldview. Descola's sharpest move was to put us inside the table too, as only one of four cells. We are no longer the destination, merely an option. An Andean offering drink to a mountain and you measuring soil acidity in a laboratory use two internally organized ways of assembling a world. This does not mean science is ineffective. It means the intuition that there is only one correct way to approach the world is itself characteristic of one cell in the table.

Now shaman. The word's origin is clear: it comes from shaman, commonly transliterated with an initial sh sound, in Siberian Tungusic languages such as Evenki. Originally it denoted a particular regional specialist who negotiated with spirits to heal, divine, and mediate between human and divine worlds. In the mid-twentieth century, scholar of religion Mircea Eliade published Shamanism: Archaic Techniques of Ecstasy, extending the term into a global type defined by ecstasy: journeys of the soul outside the body, ascending and descending between worlds. This influential extension also caused problems. Today, from Siberia to the Amazon, from Mongolia's grasslands to television documentaries, any combination of drumming, trance, and spirit communication prompts the label "shaman." Yet those labeled often do not use the word themselves, and their roles, lineages, training, and prohibitions vary enormously. Calling them all shamans resembles grouping surgeons, counselors, village officials, and ministers as "white coats" because all handle human troubles.

The consumer end is worse. In the 1980s, anthropologist Michael Harner distilled "core shamanism," stripping diverse practices of their specific cultures and retaining a general recipe of drumming, journeying, and finding power animals, then teaching courses. Shaman became a bestselling label in the spiritual marketplace: weekend workshops, paid soul journeys, and tourist shaman experiences. A circuit was complete: a word taken from Siberia $\rightarrow$ spread by scholarship into a global label $\rightarrow$ made into a commodity by markets $\rightarrow$ sold to curious urban consumers. People in its place of origin received no share and often did not agree.

Together, these word histories offer the same lesson: power and markets stand between a category's birthplace and its eventual range. Asking who coined a word, what it originally meant, and who profits from it equips you with something more durable than any individual fact.

\section{Today, Outside the Museum}
This is not merely the past of history lessons. Some bills for misreading are still being paid; others are only beginning to be settled.

First consider the cost. In the late nineteenth century, the Ghost Dance movement arose on North America's Great Plains. The Paiute prophet Wovoka taught that through dancing, singing, and peaceful living, deceased ancestors and buffalo would return and white oppression would pass. This was a peaceful movement of end-times hope that forbade violence. Yet as its name and images passed through successive accounts to white officials and newspapers, it became "Indians performing a war dance and preparing rebellion." In December 1890, at Wounded Knee in South Dakota, the U.S. Army surrounded a group of Ghost-Dancing Lakota on their way to surrender. After gunfire began, more than two hundred people died, many women and children. A misread ritual was paid for in human lives. Lakota people still commemorate Wounded Knee each year; wounds left by misreading heal on a scale of centuries.

Now the beginnings of redress. In 1990, the United States passed legislation requiring federally funded museums to return Indigenous human remains and sacred objects to their communities. Objects like the masks and drums in display cases increasingly began their journeys home. Many communities first put returned sacred objects back to work: these are not exhibits but living points of relationship. In 1992, Australia's High Court, in the Mabo decision, rejected terra nullius, the legal fiction that Australia belonged to nobody before Europeans arrived. Supporting that fiction for two centuries was the old judgment that its people had no real religion, law, or land system. In 2019, Australia permanently prohibited visitors from climbing Uluru because it is sacred to Anangu people. For decades, "It's just a rock" and "This is our church" had confronted one another on the same rock; finally, law sided with the latter for the first time. In 2016, protests against a pipeline project at Standing Rock Sioux lands brought "water is alive" into global news. Many people first recognized that, for its speakers, this was not an environmental slogan but an entire worldview of kinship between humans and rivers.

The echoes reach closer too. Some books still place primitive religion and totem worship on the lowest rung of religious evolution. Tourism films package Evenki people, shamans, and sacred drums as exotic scenery. Social-media quizzes about your power animal turn another community's sacred relationships into traffic-generating toys. This may not be malicious, but malice is unnecessary: a crooked ruler and a copying machine are enough to make misreading reproduce itself.

You can now recognize the common structure: a term emptied of context, an account stripped of its five checkpoints, a living tradition treated as past tense. The equipment is in your hands. The question is whether you will use it.

\section[Who May Judge Complexity?]{Your Question: Does Complexity Need an Admission Ticket?}
Return to the museum basement and the mask looking at you through glass.

You now know that it came from a hot gathering house where a whole cosmos was transmitted through song and bodies, and someone could hear a single mistaken word. It acquired the label "shamanic ceremonial object," a label originating in Siberia whose fitness for use is already questionable. None of its owners signed off at the checkpoints of collection, transport, cataloging, or exhibition. Its original owners' descendants are alive today, still singing, still pursuing legal claims for their belongings, and still insisting that theirs is a living tradition, not a fossil of someone else's childhood.

The chapter's question is now fully yours:

\textbf{By what authority, and with what ruler, do we measure the complexity of a tradition without scriptures?}

Before answering, review the accounts. One side says comparison needs standards and teaching needs simplification. If every tradition is placed beyond classification, speech itself becomes difficult; this chapter has used religion, belief, and ritual throughout. The other says rulers are human-made. Their makers have histories and their users have interests. For more than a century, calling something simple has prepared the way for conquest and forgetting. The gunfire at Wounded Knee began with the misjudgment that people were performing a war dance.

Give your answer and reasons. Should the mask remain behind museum glass as humanity's shared heritage, available for everyone to view? Or should it return to its community, even if that removes it from your sight or means burying it again in the tundra? Does your word "complex" describe its qualities, or announce your license to judge them?

There is no standard answer. But notice your few seconds of hesitation, when you begin to feel that this is not a question answered with an effortless "Of course we should..." Those seconds are what the chapter really hopes to leave you. Their name is respect.
